OCL2Cypher: A Semantics-Preserving Transformation
from UML/OCL Models to Property Graphs
Quoc-Bao Hoang, Hoang-Viet Trinh La, and Duc-Hanh Dang⋆
Faculty of Information Technology, VNU University of Engineering and Technology
VNU Town, Hoa Lac, Hanoi, Vietnam
{23020012 | 24025181 | hanhdd}@vnu.edu.vn
Abstract. UML models are commonly constrained using the Object Constraint
Language (OCL), but validating such constraints over a property-graph representation requires more than generating executable graph queries: the generated
query must return exactly the same set of violating context objects as source
OCL evaluation. Existing work has studied UML-to-graph representations and
OCL-to-query compilation, yet it often lacks an end-to-end semantic guarantee
that jointly relates source validation semantics, graph representation, the generated query, and stable-ID projection. We present OCL2Cypher, a representationaware method that translates a UML snapshot to a property graph and invariants
to graph queries, through typed Core and Q forms compiled to Cypher queries
under a shared representation specification Σ. We ensure its correctness obligations explicitly by checking properties during the transformation: representation
adequacy, typed translation, violation projection, and graph query realization and
execution. Under the stated obligations, their composition conditionally preserves
the exact violation-ID set for finite, well-formed snapshots and invariants in the
supported fragment. A subset of the formal obligations is checked in Lean. Prototype evaluation on the graph database Neo4j yields 199 consistent intermediate
comparisons, 23 negative cases outside the supported fragment or violating admission conditions that are rejected at the expected boundary, and 85 executions
whose violation-ID sets exactly match source evaluation. These experiments provide finite implementation evidence and do not extend the theorem’s scope.
Keywords: Object Constraint Language ꞏ Cypher ꞏ Property graph ꞏ Semantic
preservation ꞏ Model transformation
1 INTRODUCTION
In Model-Driven Engineering (MDE), UML defines model structure and the Object
Constraint Language (OCL) specifies invariants over model instances [16,4]. Graphbacked systems such as Neo4EMF and NeoEMF store model data outside the original
OCL evaluator [3,7]; Mogwaï compiles OCL to Gremlin [8,5], and UMLtoGraphDB
combines UML-to-graph mapping with OCL-to-Gremlin transformation [9].
This shift creates a semantic boundary: an executable graph query can still disagree
with source OCL if representation and compilation treat navigation, multiplicity, inher-
⋆ Corresponding author
itance, collections, or object identity differently. Prior work addresses OCL–SQL correctness [15], translations to graph constraints [17], and Cypher–SQL equivalence [12],
but current approaches have not yet jointly related source semantics, graph representation, generated-query semantics, and violation-ID projection. We therefore seek explicit
conditions under which generated Cypher returns exactly the source OCL violation-ID
set.
OCL2Cypher addresses this problem through two coordinated branches: one maps
a conforming UML snapshot to a property graph, and the other compiles an OCL invariant to read-only Cypher. A shared representation specification governs identifiers,
types, attributes, associations, multiplicities, and inheritance. Typed intermediate forms
expose representation-dependent decisions before query realization. We make the obligations across representation, translation, violation projection, realization, and execution explicit; unsupported constructs are rejected rather than approximated.
This paper makes two main contributions:
– Representation-aware OCL-to-Cypher compilation. We provide a typed transformation from a supported OCL fragment to read-only Cypher under the shared representation specification.
– Assurance framework. We identify and compose the obligations for representation
adequacy, semantic preservation, violation projection, and query realization and execution, yielding conditional exact preservation of violation IDs.
Section 2 gives the running example; Sects. 3–4 present and specify the method;
Sect. 5 establishes correctness; Sect. 6 evaluates the prototype; and Sects. 7–8 discuss
related work and conclude.
2 MOTIVATING EXAMPLE AND SEMANTIC BACKGROUND
This section uses a running example to expose a representation/translation mismatch
and introduce the required semantic concepts.
2.1 Motivating Example
Consider a UML schema with class Branch and two association ends: employee of type
Employee[*] and manager of type Employee[0..1]. Class Employee has attribute
salary: Integer. The snapshot in Fig. 1 contains a Branch with stable identifier
oneBranch and two Employee objects, Alice and Bob. Both are employees, and Alice
is also the manager; their salaries are 50 and 70, respectively.
The following invariant requires the manager’s salary to exceed that of every other
employee in the same branch:
context Branch inv ManagerSalaryDominates:
self.employee->forAll(e |
e <> self.manager implies
self.manager.salary > e.salary)
For Alice, the antecedent of the implication is false. For Bob, the antecedent is true,
but 50 > 70 is false. Hence, forAll evaluates to false and the source validation observation is {oneBranch}.
oneBranch : Branch
Alice : Employee
salary = 50
Bob : Employee
salary = 70
employee
manager
(a) An example UML snapshot
Transformation
:Branch
id=oneBranch
:Employee
name=Alice
salary=50
:Employee
name=Bob
salary=70
MANAGER
EMPLOYEE EMPLOYEE
(b) Property-graph representation
of the example snapshot
Fig. 1. Example snapshot and property-graph representation. Open UML arrowheads denote navigability, not generalization; graph arrows show physical relationship direction.
In Fig. 1, EMPLOYEE is stored from Employee to Branch, whereas self.employee
navigates in the opposite logical direction:
Naive: MATCH (b:Branch)-[:EMPLOYEE]->(e:Employee)
Required: MATCH (b:Branch)<-[:EMPLOYEE]-(e:Employee)
These are partial patterns, not complete validation queries. Equating logical with
physical direction produces an empty employee collection, makes forAll true, and
misses oneBranch, although the query remains executable. Representation and compilation must therefore agree on navigation, multiplicity, inheritance, collections, and
stable-ID projection.
Semantic background. A snapshot supplies the objects, values, and links conforming
to a UML schema. OCL evaluates an invariant for every context-class object [16]; we
preserve their violation-ID set, not one aggregate Boolean. Property-graph relationships
are directed [1,2], whereas OCL navigation follows the association end and receiver
type. The fragment also distinguishes an empty collection from typed bottom: forAll
on an empty collection is true, and any non-true invariant value is violating.
3 METHOD OVERVIEW
OCL2Cypher translates an OCL invariant under a shared representation specification Σ
and assumes a graph G satisfying AdequateΣ(M, G). Figure 2 separates this transformation from the prototype graph builder and execution context, which are not formally
verified.
Core is a source-oriented typed intermediate representation of the resolved and normalized OCL computation. Q is a representation-aware typed query IR that makes graph
access, inheritance, collection processing, and three-valued iterator decisions explicit
before Cypher realization.
Write the representation contract schematically as Σ = (Σid, Σtype, Σattr, Σassoc, Σinh,
Σcoll). The graph branch is schematically M
BuildΣ −−−−→ G: each object becomes an observable node with its stable ID and type, attributes become properties, and links follow
UML
schema
OCL
invariant
Frontend Core Q
Cypher
realization
Generated
Cypher query
Σ
Transformation pipeline
Snapshot
M
Graph
construction
Property
graph
G
Execute
Cypher on G
Target
violation IDs
Evaluate OCL
on M
Source
violation IDs
Compare
violation IDs
input/output
IR
process
shared spec.
data flow
Σ dependence
Fig. 2. OCL2Cypher transformation and observation workflow.
Σassoc. The other components govern class keys, inheritance-aware extents, multiplicity,
collection kind, and duplicates. Section 5.1 states the builder’s required output contract
as AdequateΣ(M, G).
FrontS resolves and normalizes the invariant to Core; TΣ lowers Core to Q, exposing graph-access, collection, and inheritance decisions; and RealizeΣ realizes the
Violations(K, q) wrapper as read-only Cypher. Unsupported constructs are rejected.
4 SPECIFICATION OF THE OCL2CYPHER
TRANSFORMATION
This section specifies three boundaries: OCL → Core performs source analysis, Core
→ Q exposes representation decisions, and Q → Cypher realizes the query. Snapshot
representation and query transformation share Σ.
4.1 OCL-to-Core Frontend
Given a UML schema S, FrontS is syntax-directed over the supported fragment: it resolves keys and binders, checks types, determines multiplicity and collection information, normalizes admitted constructs to Core, and rejects the rest. It returns a contextclass key K and expression e. Let Σ govern subsequent lowering and realization. A
variable environment maps resolved binders, including self, to typed values; ρ[x 7→ v]
is functional update. A successful transformation satisfies
FrontS(inv) = Success(K, e),
TΣ(e) = Success(q),
RealizeΣ(Violations(K, q)) = Success(s, ∆, πgen).
(1)
Here e ∈ ExprCore[Boolean], q ∈ ExprQ[Boolean], and (s, ∆, πgen) is the Cypher artifact. These stages use Σ, not G; Violations carries K for projection, and failure is not
an empty result.
After name resolution and multiplicity analysis, the running example has the following Core AST shape (resolved keys and κ are explicit):
ForAll(
NavManyκ
(self, remployee), x,
Bin⇒
(
Bin̸=(x, NavOne(self, rmanager)),
Bin>(Attr(NavOne(self, rmanager), asalary), Attr(x, asalary)))).
Here remployee, rmanager, and asalary are resolved schema keys. Source multiplicity distinguishes to-one from to-many navigation; the source collection type fixes κ, subsequently
checked against Σ by TΣ.
Scope. The fragment admits finite well-formed snapshots; three-valued Boolean, String,
object, Integer, and Real values; inheritance-aware navigation; finite flat Sets and Bags;
and select, scalar-body collect, forAll, exists, isEmpty, notEmpty, size, and
sum. It excludes ordered or nested collections, Tuples, recursion, order-sensitive operators, surface null/invalid, and any. Integers must remain in INT64 and Reals exactly
representable; unchecked runtime range and precision conditions are P5 premises of
Theorem 2.
Typed bottom represents an undefined partial computation. For defined-value domain Dτ , let
D⊥
τ
:= Dτ ∪ {⊥τ }, D⊥
B := {T, F, ⊥B}.
Here B is Boolean. Equality is typed and total: objects compare by class key and stable
ID, same-typed bottoms are equal, and the result is never ⊥B. Also ¬⊥B = ⊥B, ∧ is
F-dominant, ∨ is T-dominant, and β1 ⇒ β2 := ¬β1 ∨ β2.
4.2 Core-to-Q Translation
Abstract syntax. Core and Q are modeled as mathematical typed ASTs. For L ∈
{Core, Q}, ExprL and PlanL denote the expression and collection-plan AST sorts,
respectively, and ExprL[τ ] contains well-typed expressions of type τ . This paper-level
calculus abstracts from both the production Java classes and the Lean datatypes; it is
neither a concrete textual grammar nor a metamodel. Variables eL, PL range over terms,
with collection processing factored into plans. The two IRs share these AST constructor
families:
eL ::= cτ | ⊥τ | x | Let(x, eL, eL) | If(eL, eL, eL)
| Attr(eL, a) | NavOne(eL, r) | Uno1
(eL) | Bino2
(eL, eL)
| Litκ(¯eL) | Materialize(PL),
PL ::= NavManyκ
(eL, r) | Collect(PL, x, eL) | Distinct(PL).
The language-specific forms are
ExprCore ∋ ForAll(PC , x, eC ) | Exists(PC , x, eC ),
ExprQ ∋ ForAll3(PQ, x, eQ) | Exists3(PQ, x, eQ),
PlanCore ∋ From(eC ) | ScanK | Select(PC , x, eC ) | LetP(x, eC , PC ),
PlanQ ∋ FromCollection(eQ) | ScanClassK | Filter(PQ, x, eQ) | PlanLet(x, eQ, PQ).
Here κ ∈ {Set, Bag}; U contains Boolean negation, numeric unary operations, and
collection observations; and O2 contains Boolean, typed comparison/equality, and numeric binary operations. Thus o1 ∈ U, o2 ∈ O2; a, r, K are resolved keys, bars are
finite sequences, and displayed syntax omits type/Σ data. Violations(K, q) is a top-level
wrapper outside both AST sorts.
The evaluations [[·]]O, [[·]]Core, and [[·]]Q are structural; Core/Q expression and plan
evaluation is mutually defined into typed-bottom domains. The artifact supplies every
constructor clause.
TΣ consumes typed Core and Σ, without observing G. Define successful translation
by
Σ ` t ⇝ t
′
:⇐⇒ TΣ(t) = Success(t
′
).
On admitted types, TyΣ : TypeCore → TypeQ is identity on primitives/class keys
and homomorphic on Set and Bag; extend it pointwise to typing contexts. Translation
preserves sort (expression to expression, plan to plan) and typing:
Γ `Core t : τ ∧ Σ ` t ⇝ t
′ =⇒ TyΣ(Γ) `Q t
′
: TyΣ(τ ).
Table 1. Core-to-Q forms and semantic preservation obligations.
Core form Q form Applicable semantic preservation obligation
Constants, variables,
bottom, Let/If/Attr
cτ , x, ⊥τ , Let/If/Attr Corresponding values, bottoms, branches, and attribute
observations agree.
NavOne(e, r) NavOne(q, r) Source navigation agrees with the encoding, direction, target,
and multiplicity in Σ(r).
NavManyκ(e, r) NavManyκ(q, r) Source occurrences agree with Σ(r); preserve Set support or
each Bag occurrence.
Operations; collection
observations
Corresponding primitives Admitted primitive meanings agree.
Lit; Select; Collect;
Distinct
Lit; Filter; Collect;
Distinct
Selection preserves retained occurrences and collection kind;
collect preserves occurrence multiplicity as a Bag; Distinct
preserves Set support.
From; Materialize; LetP FromCollection;
Materialize; PlanLet
Preserve collection observations, occurrence multiplicity, and
whole-bottom behavior across expression/plan boundaries.
forAll/exists ForAll3/Exists3 Finite source; distinguish empty collections from whole-bottom.
ScanK; invariant body ScanClassK;
Violations(K, q)
[[ScanK]]Core = InstM(K) and
[[ScanClassK]]Q = ExtentΣ
G(K); stable-ID projection retains
exactly the cases where q ̸= T.
Translation establishes key/binder identity, sort preservation, and typing. Table 1 exhaustively covers supported constructor families and lists the remaining semantic obligations. A structurally non-trivial rule is
Σ ⊢ PC ⇝ qP Σ ⊢ b ⇝ b
′
Σ ⊢ ForAll(PC , x, b) ⇝ ForAll3(qP , x, b′)
.
The rule retains binder scope; exists similarly uses Exists3. For plan result C ∈
D⊥
Set(τ)
or C ∈ D⊥
Bag(τ)
, let Bb(v) := [[b
′
]]Q(G, η[x 7→ v]) for v ∈ Occ(C). Wholebottom is distinct from the empty collection; ForAll3 follows the ordered cases below:



⊥B, C is whole-bottom,
F, ∃v ∈ Occ(C) : Bb(v) = F,
⊥B, ∄v ∈ Occ(C) : Bb(v) = F ∧ ∃v ∈ Occ(C) : Bb(v) = ⊥B,
T, otherwise.
Here Occ(C) enumerates Bag occurrences or Set support; empty input yields T. Exists3
is dual, prioritizing T, then ⊥B, then F; empty input yields F.
4.3 Q-to-Cypher Realization
In (1), ∆ = (Dompub, idCol, decode) gives the caller-parameter domain, returned ID
column, and decoder to I ⊆ ID. Parameters πpublic and realization-generated πgen must
have disjoint names.
The emitted read-only fragment uses MATCH patterns with WHERE predicates, WITH,
UNWIND, RETURN, CASE, comprehensions, quantified predicates, and EXISTS/COLLECT
subqueries. Writes, procedures, dynamic labels, and OPTIONAL MATCH are excluded.
Table 2. Q-to-Cypher realization patterns.
Q form or wrapper Representative Cypher realization pattern
ScanClassK COLLECT subquery with a type- and inheritance-aware context scan determined by Σ.
cτ , ⊥τ , x; Let/If/Attr;
Un/Bin
Tagged constants and bottoms, bound variables, bind-once comprehensions, CASE, checked
property access, and typed primitive expressions with explicit bottom propagation.
NavOne(q, r) Σ-directed MATCH; exactly one collected target yields the corresponding defined target value,
and any other cardinality yields typed bottom.
NavManyκ(q, r) Σ-directed matching in COLLECT; DISTINCT for Set, with occurrences retained for Bag.
Filter(P, x, b) Bottom-aware comprehension retaining occurrences satisfying b.
Collect(P, x, e) Occurrence-preserving comprehension yielding a Bag.
Litκ; Materialize;
FromCollection; Distinct;
PlanLet
Tagged collection values; expression/plan bridges retain items and whole-bottom; Set
realization removes duplicate atomic values, and plan binding uses bind-once realization.
ForAll3/Exists3 Explicit defined/bottom cases, quantified checks, and CASE implement the three-valued
priority rules.
Violations(K, q) Context scan retaining non-true results; RETURN DISTINCT projects stable IDs.
Table 2 covers every supported Q family and the Violations wrapper at the level of
realization patterns, not one-to-one syntax; the artifact contains the production generator.
5 CORRECTNESS OF THE OCL2CYPHER
TRANSFORMATION
The correctness argument uses a representation-adequacy contract and factors compilation through semantic preservation from OCL to Core and from Core to Q, followed by
violation projection and graph query realization and execution.
5.1 Formal Setup and Observations
Let a UML schema S, finite snapshot M, graph G, and representation specification Σ
be given. Conforms(M, S) requires declared object types, typed attributes and association endpoints, valid qualifiers, and inheritance-consistent membership. WFS(M) additionally requires injective stable IDs, no dangling endpoints or ambiguous scalar slots,
and deterministic upper-one observations, with absence as typed bottom. Here “stable”
means preserved by this representation, not across model versions.
ObjObsΣ(G) contains object-representing graph nodes, excluding schema, metadata, and reification nodes. Source and graph type functions share a class-key domain,
while idM : Obj(M) → ID and idG : ObjObsΣ(G) → ID share the codomain ID, with
idM injective. A bijection γ : Obj(M) ∼= ObjObsΣ(G) pairs each source object with its
graph representation. Its type-directed lifting fixes primitives and typed bottom, maps
objects by γ, and maps collections pointwise: γb(c) = c, γb(⊥τ ) = ⊥τ , γb(o) = γ(o),
and γb(κhviiI ) = κhγb(vi)iI . Here I indexes Set support or individual Bag occurrences,
preserving Bag multiplicity.
For fixed Σ, slotG, NavΣ
G, and ExtentΣ
G observe graph attributes, logical navigation,
and inheritance-aware extents; their source counterparts are slotM, NavM, and InstM.
Definition 1 (Observational representation adequacy). We say that G is Σ-adequate
for M, written AdequateΣ(M, G), iff there is a bijection γ : Obj(M) ∼= ObjObsΣ(G)
such that, for every object o, every attribute a and association end r applicable to o,
and every class key K in the supported domain,
idG(γ(o)) = idM(o), typeG(γ(o)) = typeM(o),
γb(slotM(o, a)) = slotG(γ(o), a),
γb(NavM(o, r)) = NavΣ
G(γ(o), r), γ[InstM(K)] = ExtentΣ
G(K).
Whenever adequacy is assumed, fix its witness γ and lifting γb, and write
v 'Σ w :⇐⇒ w = γb(v),
ρ 'Σ η :⇐⇒ dom(ρ) = dom(η) ∧ ∀x ∈ dom(ρ), ρ(x) 'Σ η(x).
The theorem assumes adequacy; Sect. 6 evaluates the non-formalized graph builder
empirically.
Let invariant inv have body φ and resolved context key K, and let q be its Q body.
Source, Q, and Cypher observations are
ViolO(M, inv) := {idM(o) | o ∈ InstM(K), [[φ]]O(M, o) 6= T},
ObsQ(G, K, q) := [[Violations(K, q)]]Q(G),
ViolC (G, s; π, ∆) = I ⇐⇒ [[s]]C (G; π, ∆) = Success(I).
ViolC is a partial function, defined only when parsing, execution, and decoding succeed;
failure is not interpreted as the empty set. The Q wrapper scans the extent, retains results
other than T, and projects stable IDs.
5.2 Preservation across OCL–Core and Core–Q
Lemma 1 (Frontend preservation). Under Conforms(M, S) ∧ WFS(M), let source
invariant inv belong to the supported fragment from Sect. 4.1 and be well-typed with
respect to S. Let its body be φ and its resolved context-class key be K. If FrontS(inv) =
Success(K, e), then for every o ∈ InstM(K),
[[φ]]O(M, o) = [[e]]Core(M, ρo),
where ρo maps self to o.
Proof sketch. Resolution preserves referenced objects and binder scope; normalization
replaces admitted surface forms by equivalent Core constructors, while type checking
excludes inadmissible cases.
Theorem 1 (Core-to-Q preservation). Assume AdequateΣ(M, G) and ρ 'Σ η. Let
e ∈ ExprCore[τ ] and q ∈ ExprQ[TyΣ(τ )] such that Σ ` e ⇝ q. If the applicable
semantic preservation obligations in Table 1 hold, then
[[e]]Core(M, ρ) 'Σ [[q]]Q(G, η).
This expression result is proved mutually with the corresponding preservation result for
PC ∈ PlanCore and PQ ∈ PlanQ.
Proof sketch. Simultaneous induction over expression and plan translation handles literals, variables, and bottom directly; attributes and navigation use adequacy; collection
cases use preservation of Set support and Bag occurrences; and iterators apply their ordered three-valued cases. Primitive cases use the applicable obligations. The technical
report gives the complete proof.
5.3 Violation Projection and Graph Query Realization and Execution
Lemma 2 (Violation-projection preservation). Assume AdequateΣ(M, G) and
∀o ∈ InstM(K) : [[e]]Core(M, ρo) 'Σ [[q]]Q(G, ηγ(o)).
Here ρo and ηγ(o) map self to o and γ(o), respectively. Then
{idM(o) | o ∈ InstM(K), [[e]]Core(M, ρo) 6= T} = ObsQ(G, K, q).
Proof. Adequacy pairs the extents and stable IDs, while Boolean preservation retains
the condition of being different from T. The wrapper removes duplicate IDs in the final
observation.
Let TargetΣ contain typed, read-only realization records. Each p ∈ TargetΣ carries a target Boolean body, ID-projection predicate, and serialization metadata, interpreted as B
p
T
and Retp
; it is an abstract realization IR, not a Cypher AST. The relation
WitnessΣ associates p, G, K, q with (s, ∆, πgen) exactly when p is a well-typed realization of Violations(K, q) under Σ and serializes to that artifact. It records provenance
and artifact correspondence, not R1–R3; G interprets p, but realization does not read G.
Define
E := ExtentΣ
G(K), BQ(g) := [[q]]Q(G, ηg),
where ηg(self) = g. These observations have types
E ⊆ ObjObsΣ(G),
BQ, Bp
T
: E → D⊥
B, Retp
: ID → Prop.
R1–R2 concern realization. Given successful execution and decoding [[s]]C (G; π, ∆) =
Success(I) with I ⊆ ID, R3 is the execution-and-decoding obligation:
R1 : ∀g ∈ E, BQ(g) = B
p
T
(g),
R2 : ∀j ∈ ID, Retp
(j) ⇐⇒ ∃g ∈ E : idG(g) = j ∧ B
p
T
(g) 6= T,
R3 : ∀j ∈ ID, j ∈ I ⇐⇒ Retp
(j).
Lemma 3 (Generated-query refinement). Assume
RealizeΣ(Violations(K, q)) = Success(s, ∆, πgen).
Suppose there is a p such that WitnessΣ(p; G, K, q, s, ∆, πgen) and R1–R2 hold for p.
Let πpublic ∈ Dompub satisfy
dom(πpublic) ∩ dom(πgen) = ∅, π = πpublic ] πgen.
If [[s]]C (G; π, ∆) = Success(I) and R3 holds for p, I, then
I = ObsQ(G, K, q).
Proof sketch. R1 preserves body values, R2 selects exactly the stable IDs whose values
are non-true, and R3 gives the same membership condition for I.
5.4 End-to-End Theorem and Mechanized Evidence
Write CompileS,Σ(inv) = Success(s, ∆, πgen) when some K, e, q satisfy (1); this merely
names the composition.
Theorem 2 (Conditional exact preservation of the violation-ID set). Let S, M, G, Σ,
invariant inv, and compilation result (s, ∆, πgen) be given. Assume:
P1. Source admissibility: Conforms(M, S)∧WFS(M); M is finite, and inv is supported
and well-typed under S.
P2. Representation adequacy: AdequateΣ(M, G).
P3. Compilation: CompileS,Σ(inv) = Success(s, ∆, πgen). Take K, e, q from Eq. (1).
P4. Refinement: There is a p such that WitnessΣ(p; G, K, q, s, ∆, πgen); all applicable
Table 1 obligations and R1–R2 hold.
P5. Runtime: πpublic ∈ Dompub, dom(πpublic) ∩ dom(πgen) = ∅, and π = πpublic ] πgen.
Applicable runtime numeric obligations hold; [[s]]C (G; π, ∆) = Success(I) with
I ⊆ ID. R3 holds for the same p, I.
Then
ViolO(M, inv) = ViolC (G, s; π, ∆).
Proof sketch. P1 and Lemma 1 preserve source evaluation in Core. P2–P5 supply adequacy, translation, preservation, and numeric premises; Theorem 1 and Lemma 2, followed by Lemma 3, yield
ViolO(M, inv) = ObsQ(G, K, q) = ViolC (G, s; π, ∆).
Thus the obligations are sufficient and compositional; the theorem does not assert that
production runs automatically satisfy adequacy, R1–R3, or runtime premises.
Scope of mechanization. Lean’s trusted kernel checks the formal normalization, Coreto-Q and Q-to-target preservation, and their composition, without unchecked placeholders or custom axioms. Adequacy, production compiler/graph builder, Cypher execution,
decoding, and runtime numeric assumptions remain outside this scope and are evaluated
separately in Sect. 6.
6 EVALUATION
Prototype implementation. The prototype builds a snapshot graph and lowers an invariant through Core and Q to read-only Cypher. Both paths use Σ and reject inadmissible
inputs. The artifact1 provides code, corpus, results, reproduction instructions, and separate adequacy checks (excluded from RQ counts) for IDs, types, attributes, navigation,
and extents.
Research questions and evaluation design. We evaluate OCL2Cypher using three research questions.
RQ1—Core–Q agreement. On prototype-admitted inputs, do Core and Q produce the
same values, collection observations, and violation-ID sets?
RQ2—Rejection. Are inadmissible inputs rejected at the prescribed transformation
boundary?
RQ3—End-to-end execution. Does executing the generated Cypher query on Neo4j
return the same violation-ID set as source OCL evaluation?
Units are invariant–snapshot pairs. RQ1 compares Core/Q values and violation IDs;
RQ3 executes generated Cypher on Neo4j. RQ2 contains 23 negatives: eight require an
exact stage/code and fifteen require rejection before Core. The source oracle evaluates
frontend-produced Core directly on the snapshot. It shares data representations and frontend/lowering code, but bypasses Q, graph building, realization, Neo4j, and decoding;
fixed ID fixtures provide an additional check. Thus RQ3 spans boundaries including R3
but does not establish R3 independently.
The feature matrix also tests prototype extensions—reject, Set operations,
allInstances, and type tests/casts—outside the theorem profile. These tests do
not enlarge the theorem’s supported language.
Table 3. RQ results: a/b successful checks over total; dash means inapplicable.
Subject Focus RQ1 RQ2 RQ3
Library Collections, bottom, and numeric bounds 78/78 1/1 –
CarRental Domain constraints; frontend admission 19/19 6/6 19/19
DeepHierarchy Inheritance-aware context extent 36/36 – –
Medical Navigation and iterators 36/36 – 36/36
Feature matrix Profile risks and implementation extensions 30/30 1/1 30/30
Boundary corpus Out-of-scope inputs – 15/15 –
Total per unit 199/199 23/23 85/85
Subjects and results. Table 3 shows that RQ1 passed all 199 intermediate comparisons;
RQ2 rejected all 23 designated negatives—8/8 at their exact expected stage/code and
15/15 before Core as required; and RQ3 matched the oracle in all 85 executions, with
no missing or spurious IDs, using Cypher 5 on Neo4j 2026.06.0 Enterprise. These finite
results cover only the reported corpus and configuration; they imply neither general
correctness nor performance or scalability.
1
https://github.com/bao2811/OCL2Gratra.
Threats to validity. Shared frontend/lowering code and data can correlate prototype and
oracle faults. We mitigate this by evaluating Core directly while bypassing Q, graph construction, realization, Neo4j, and decoding, and by checking fixed ID fixtures. The corpus targets distinct risks and rejection cases but not all OCL or industrial-scale models;
RQ3 uses one Neo4j configuration. Adequacy checks and formal obligations provide
complementary evidence but do not establish portability or correctness probabilities for
unseen inputs.
7 RELATED WORK
For OCL execution, USE evaluates OCL directly [11], while SQL-PL4OCL and
OCL2PSQL compile it to SQL [10,14]; neither line connects a graph representation to
its validation observation. For equivalence, Nguyen and Clavel [15] give an SMT-based
method for proving manually written SQL implementations correct for a significant
OCL subset. Radke et al. prove satisfaction equivalence between Essential OCL invariants and nested graph constraints [17]. Graphiti checks Cypher–SQL equivalence for
supported query fragments modulo a supplied database transformer [12]; none covers
the complete UML/OCL-to-Cypher validation pipeline.
For graph-backed validation, Neo4EMF and NeoEMF [3,7] provide persistence.
Prior work [5,8] has compiled OCL into Gremlin traversals through Mogwaï; GremlinATL [6] translates ATL transformations, whose expressions use OCL, into Gremlin
scripts. Other approaches [9,13] map UML schemas or class diagrams to graphdatabase representations and labeled property graphs. In contrast, this work specifies
a representation-aware translation of OCL invariants into Cypher and identifies the
semantic preservation obligations required for a conditional end-to-end correctness
result.
8 CONCLUSION
This paper introduced OCL2Cypher, a representation-aware method for compiling supported OCL invariants into read-only Cypher queries over property graphs. Its two
main contributions are a typed OCL–Core–Q–Cypher transformation under a shared
representation specification and an assurance framework that makes explicit the obligations for representation adequacy, semantic preservation, violation projection, and graph
query realization and execution. Under the stated premises, these obligations compose
to conditionally preserve the exact violation-ID set. Lean mechanically checks normalization, Core-to-Q preservation, Q-to-target preservation, and their composition, while
representation adequacy, production components, Neo4j execution, decoding, and runtime numeric assumptions remain outside the mechanized scope. Future work will extend the supported fragment to nested and ordered collections and Tuple values, and
evaluate OCL2Cypher on larger industrial models.
References
1. Angles, R., Arenas, M., Barceló, P., Hogan, A., Reutter, J., Vrgoč, D.: Foundations of modern query languages for graph
databases. ACM Comput. Surv. 50(5), 68:1–68:40 (2017). doi: 10.1145/3104031
2. Angles, R., Bonifati, A., Dumbrava, S., Fletcher, G., Green, A., Hidders, J., Li, B., Libkin, L., Marsault, V., Martens,
W., Murlak, F., Plantikow, S., Savković, O., Schmidt, M., Sequeda, J., Staworko, S., Tomaszuk, D., Voigt, H., Vrgoč,
D., Wu, M., Živković, D.: PG-Schema: Schemas for property graphs. Proc. ACM Manag. Data 1(2), 1–25 (Jun 2023).
doi: 10.1145/3589778, Art. no. 198
3. Benelallam, A., Gómez, A., Sunyé, G., Tisi, M., Launay, D.: Neo4EMF: A scalable persistence layer for EMF models.
In: Proc. 10th Eur. Conf. Model. Foundations Appl. Lect. Notes Comput. Sci., vol. 8569, pp. 230–241 (2014). doi: 10.
1007/978-3-319-09195-2_15
4. Cabot, J., Gogolla, M.: Object constraint language (OCL): A definitive guide. In: Formal Methods for Model-Driven
Engineering, Lect. Notes Comput. Sci., vol. 7320, pp. 58–90. Springer (2012). doi: 10.1007/978-3-642-30982-3_3
5. Daniel, G.: Efficient validation of large models using the Mogwaï tool. In: Proc. 18th Int. Workshop on OCL and Textual
Model. (OCL). CEUR Workshop Proc., vol. 2245, pp. 187–193 (2018), https://ceur-ws.org/Vol-2245/ocl_
paper_8.pdf
6. Daniel, G., Jouault, F., Sunyé, G., Cabot, J.: Gremlin-ATL: A scalable model transformation framework. In: Proc. 32nd
IEEE/ACM Int. Conf. Automated Softw. Eng. (ASE). pp. 462–472. IEEE (2017). doi: 10.1109/ASE.2017.8115658
7. Daniel, G., Sunyé, G., Benelallam, A., Tisi, M., Vernageau, Y., Gómez, A., Cabot, J.: NeoEMF: A multi-database model
persistence framework for very large models. Sci. Comput. Program. 149, 9–14 (2017). doi: 10.1016/j.scico.2017.08.002
8. Daniel, G., Sunyé, G., Cabot, J.: Mogwaï: A framework to handle complex queries on large models. In: Proc. 10th Int.
Conf. Research Challenges Inf. Sci. (RCIS). pp. 1–12. IEEE (2016). doi: 10.1109/RCIS.2016.7549343
9. Daniel, G., Sunyé, G., Cabot, J.: UMLtoGraphDB: Mapping conceptual schemas to graph databases. In: Proc. 35th
Int. Conf. Conceptual Model. (ER). Lect. Notes Comput. Sci., vol. 9974, pp. 430–444. Springer (2016). doi: 10.1007/
978-3-319-46397-1_33
10. Egea, M., Dania, C.: SQL-PL4OCL: An automatic code generator from OCL to SQL procedural language. Softw. Syst.
Model. 18(1), 769–791 (2019). doi: 10.1007/s10270-017-0597-6
11. Gogolla, M., Büttner, F., Richters, M.: USE: A UML-based specification environment for validating UML and OCL.
Sci. Comput. Program. 69(1–3), 27–34 (2007). doi: 10.1016/j.scico.2007.01.013
12. He, Y., Fang, R., Dillig, I., Wang, Y.: Graphiti: Bridging graph and relational database queries. Proc. ACM Program.
Lang. 9(PLDI), 1–25 (Jun 2025). doi: 10.1145/3729319, article 216
13. León, A., Santos, M.Y., García, A., Casamayor, J.C., Pastor, O.: Model-to-model transformation: From UML class diagrams to labeled property graphs. Bus. Inf. Syst. Eng. 66(1), 85–110 (2024). doi: 10.1007/s12599-023-00824-9
14. Nguyen Phuoc Bao, H., Clavel, M.: OCL2PSQL: An OCL-to-SQL code-generator for model-driven engineering. In:
Proc. Int. Conf. Future Data and Security Eng. (FDSE). Lect. Notes Comput. Sci., vol. 11814, pp. 185–203 (2019).
doi: 10.1007/978-3-030-35653-8_13
15. Nguyen Phuoc Bao, H., Clavel, M.: Proving correctness for SQL implementations of OCL constraints (2024). doi: 10.
48550/arXiv.2403.18599, arXiv:2403.18599 [cs.DB]
16. Object Management Group: Object constraint language, version 2.4. Tech. Rep. formal/2014-02-03, Object Management
Group (Feb 2014), https://www.omg.org/spec/OCL/2.4
17. Radke, H., Arendt, T., Becker, J.S., Habel, A., Taentzer, G.: Translating essential OCL invariants to nested graph constraints for generating instances of meta-models. Sci. Comput. Program. 152, 38–62 (2018). doi: 10.1016/j.scico.2017.
08.006